\section{各州比较：不同项目情景下的税务规划}

\subsection{得克萨斯州：临时销售和使用税豁免}

根据《得克萨斯州税法》第151.359条及151.3595条，得克萨斯州为符合条件的数据中心和符合条件的大型数据中心项目提供临时销售和使用税豁免。符合条件的数据中心在合资格采购中可免缴6.25\%的州销售和使用税，但仍须缴纳地方销售和使用税；免税期为10年或15年。对于符合条件的大型数据中心项目，合资格采购可在20年内同时免缴州和地方销售和使用税。

免税范围涵盖数据中心运营所必需且不可或缺的电力和设备，包括服务器、数据存储设备和网络连接设备。对于电力采购，通常需要开展主要用途研究，以区分应税和非应税用电，除非数据中心或项目符合适用的独立设施要求。

% 两类项目相同的共同要求合并展示，分别保留具体门槛。
\Needspace{0.60\textheight}
\begingroup
\ReportTableSetup
\begin{longtable}{|L{0.28\linewidth}|L{\dimexpr0.72\linewidth-4\tabcolsep-3\arrayrulewidth\relax}|}
\caption{资格要求}\label{tab:texas-eligibility}\\
\hline
\multicolumn{2}{|L{\dimexpr\linewidth-2\tabcolsep-2\arrayrulewidth\relax}|}{\textbf{共同要求（适用于两类项目）}} \\ \hline
\multicolumn{2}{|L{\dimexpr\linewidth-2\tabcolsep-2\arrayrulewidth\relax}|}{1. 持有或同时申请得克萨斯州销售或使用税许可证，或得克萨斯州直接缴税许可证；} \\ \hline
\multicolumn{2}{|L{\dimexpr\linewidth-2\tabcolsep-2\arrayrulewidth\relax}|}{2. 提交说明如何达到资本投入及就业创造要求的计划文件，并附拟议场地平面图或建筑示意图；} \\ \hline
\multicolumn{2}{|L{\dimexpr\linewidth-2\tabcolsep-2\arrayrulewidth\relax}|}{3. 位于得克萨斯州；} \\ \hline
\multicolumn{2}{|L{\dimexpr\linewidth-2\tabcolsep-2\arrayrulewidth\relax}|}{4. 专门新建或翻修，主要用于容纳服务器、相关设备及支持人员，以进行数据处理、存储和分发；} \\ \hline
\multicolumn{2}{|L{\dimexpr\linewidth-2\tabcolsep-2\arrayrulewidth\relax}|}{5. 具备不间断电源、发电机备用电源、完善的消防及防火系统，以及加强的实体安保措施，包括限制出入、视频监控和电子系统；} \\ \hline
\multicolumn{2}{|L{\dimexpr\linewidth-2\tabcolsep-2\arrayrulewidth\relax}|}{6. 不得主要供电信服务商用于提供电信服务；} \\ \hline
\multicolumn{2}{|L{\dimexpr\linewidth-2\tabcolsep-2\arrayrulewidth\relax}|}{7. 不受根据《税法》第313章限制数据中心财产评估价值的协议约束。} \\ \hline

\textbf{类型} & \textbf{具体要求} \\ \hline
符合条件的数据中心 &
1. 承诺在所在县创造至少20个合资格就业岗位；\par
2. 承诺在五年内投入或同意投入至少2亿美元资本；\par
3. 具有至少100,000平方英尺、供单一合资格使用者使用的建筑空间。
\\ \hline
符合条件的大型数据中心项目 &
1. 承诺在所在县创造至少40个合资格就业岗位；\par
2. 承诺在五年内投入或同意投入至少5亿美元资本；\par
3. 具有至少250,000平方英尺、供单一合资格使用者使用的建筑空间。
\\ \hline
\end{longtable}
\endgroup


可享受免税待遇的合资格数据中心，须具有至少100,000平方英尺、供单一合资格使用者使用的空间，在认证后五年内投入至少2亿美元资本，并创造至少20个合资格就业岗位。对于符合条件的大型数据中心项目，门槛提高至总面积至少250,000平方英尺、五年内资本投入至少5亿美元，以及至少40个合资格岗位。

实施要求明确允许数据中心设施为专门新建或翻修的设施。因此，现有矿场的改造不会仅因其原始用途而被排除在优惠范围之外。

为持续享受免税待遇，相关企业还必须遵守后续要求，包括履行适用的资本投入和合资格岗位承诺，向得克萨斯州审计长办公室提交规定报告，并完整保存免税采购记录，直至数据中心认证到期。审计长办公室会审查合规情况并开展审计；不符合要求的数据中心可能失去免税待遇，或被撤销认证。

\subsection{弗吉尼亚州：零售销售和使用税豁免与用电税}

《弗吉尼亚州法典》第58.1-609.3(18)条规定，符合条件的数据中心运营商及其租户购买或租赁合资格计算机设备和配套软件，可免缴零售销售和使用税（DCRSUT 豁免）。该豁免目前有效至2035年6月30日。如果数据中心满足特定资本投入及就业创造要求，可通过谅解备忘录（MOU）将豁免期延长至2040年或2050年。

合资格计算机设备包括服务器、路由器及数据存储设备，并明确涵盖支持免税设备运行的线缆、交换机、布线及类似物品。但与数据处理无关的一般建筑改良和固定设施，例如照明和围栏，不在豁免范围内；原本须缴纳零售销售和使用税、用于供电的燃料也不适用该豁免。

与得克萨斯州不同，一般地区的数据中心运营商及租户需要新增至少1.5亿美元资本投入。此外，数据中心须创造至少50个与数据中心运营或维护相关的就业岗位，每个岗位的薪酬至少达到当地通行年平均工资的150\%。在经济困难地区，法定门槛降低至7,000万美元资本投入，以及10个与运营或维护相关的岗位。

\ReportImage{3}{figure}{一般地区（灰色）与经济困难地区（蓝色）的分布；地名保留英文原图}

DCRSUT 豁免附有记录保存和报告要求。数据中心及其租户必须完整、准确地记录资本投入、新增岗位数量及其年平均工资、合资格计算机设备或配套软件支出，以及 DCRSUT 豁免的金额。数据中心须代表自身及租户，向弗吉尼亚州经济发展合作署（VEDP）和州税务部门提交年度进展报告及绩效目标完成情况的最终报告，并持续提交涵盖资本投入、数据中心岗位总数、年平均工资、设备支出及豁免金额的年度报告。

如果数据中心未能在规定期限内（通常为三年）达到最低资本投入、就业创造或其他绩效目标，相关企业必须补缴全部已获减免的销售和使用税，并支付利息。

自2026年7月1日起，弗吉尼亚州实施新的数据中心用电税。在为期两年的期间内，位于该州的数据中心每消耗一千瓦时电力，须缴纳0.011美元的税款。

\subsection{佐治亚州：州和地方销售和使用税豁免}

根据《佐治亚州规章汇编》第560-12-2-.117条，高科技数据中心及其客户购买和使用合资格高科技数据中心设备，可免缴州和地方销售和使用税。

除高科技数据中心所有者外，高科技数据中心客户也可享受该豁免。“高科技数据中心客户”指与数据中心所有者签订合资格合同的客户、租户、被许可人或最终用户。

免税范围涵盖用于建设、管理或维护数据中心环境的计算机设备及材料、软件、硬件或设备，包括应急备用发电机、配电单元和路由器。不动产及日常运营用电不在豁免范围内。

主要条件包括：申请免税的设备必须安装于相关高科技数据中心，或用于该数据中心；购买者须为持有免税证书的数据中心所有者或高科技数据中心客户。数据中心所有者申请免税证书时，不受最低建筑面积要求限制，但仍须满足适用的最低投资门槛和新增优质岗位要求。高科技数据中心客户的资格取决于所有者免税证书的有效性及双方的合同关系。

% 依据英文版表3翻译，金额单位保留为美元。
\begin{table}[htbp]
\ReportTableSetup
\caption{投资期内数据中心所有者申请免税证书的要求（适用于2022年5月9日或之后的申请）}
\label{tab:georgia-thresholds}
\begin{tabularx}{\linewidth}{|Y|Y|Y|}
\hline
\textbf{县人口} & \textbf{新增优质岗位平均数量} & \textbf{合资格累计支出（美元）} \\ \hline
超过50,000人 & 25 & 2.5亿 \\ \hline
30,000--50,000人 & 10 & 7,500万 \\ \hline
少于30,000人 & 5 & 2,500万 \\ \hline
\end{tabularx}
\end{table}


相应地，高科技数据中心所有者在享受免税待遇的同时，还须遵守法定年度报告、记录保存及其他要求。

\subsection{税务规划思路}

不同辖区的税收优惠在资格门槛、免税范围、有效期限及申请程序方面存在差异。因此，矿企需要结合每个项目的具体事实，评估税务处理及其影响。以下示例展示了具有不同特点的项目。

矿企 A 计划投资或改造一个大型数据中心，投资额超过5亿美元，建筑面积超过250,000平方英尺。如果项目位于得克萨斯州，可能符合大型数据中心项目的资格，使合资格设备和电力采购同时免缴州及地方销售和使用税。20年的免税期及电力采购被纳入豁免范围，使得克萨斯州在项目长期运营方面具有明显优势。

矿企 B 计划改造一个位于低人口县的矿场，投资额低于2亿美元，创造的岗位也较少。如果项目位于佐治亚州的低人口县，只需达到2,500万美元投资及5个新增优质岗位的门槛，便可能符合税收优惠条件。这一路径使项目能够根据自身特点取得税务减免，但也面临用电不享受豁免、优惠窗口较短等不足。如果设备采购延迟，剩余优惠窗口也会相应缩短。

矿企 C 计划保留园区及电力基础设施，由多家租户分别采购服务器和网络设备。如果项目位于弗吉尼亚州，该州框架允许租户通过参与安排加入免税计划，有利于更灵活的商业结构。如果项目位于经济困难地区，7,000万美元投资门槛及10个岗位的要求也可能具有吸引力。但同期征收的用电税可能削弱设备免税带来的收益。

对于转型中的矿企，免税期未来还能持续多久、用电是否享受税收优惠、租户是否纳入豁免范围、项目是否满足适用资格门槛，以及所在辖区是否已暂停或准备暂停新的免税申请，都是整体税务规划中需要考虑的重要因素。

\FloatBarrier
